% Generated draft: requires mathematical and source review.
\documentclass[11pt]{article}
\usepackage[margin=25mm]{geometry}
\usepackage{fontspec}
\setmainfont{Latin Modern Roman}
\usepackage{amsmath,amssymb,mathtools,unicode-math}
\setmathfont{Latin Modern Math}
\usepackage{enumitem,needspace,xurl,hyperref}
\setlength{\parindent}{0pt}
\setlength{\parskip}{5pt}
\setlength{\emergencystretch}{3em}
\newcommand{\N}{\mathbb{N}}
\newcommand{\Z}{\mathbb{Z}}
\newcommand{\Q}{\mathbb{Q}}
\newcommand{\R}{\mathbb{R}}
\newcommand{\C}{\mathbb{C}}
\newcommand{\F}{\mathbb{F}}
\newcommand{\A}{\mathbb{A}}
\newcommand{\PP}{\mathbb{P}}
\newcommand{\E}{\mathbb{E}}
\newcommand{\eps}{\varepsilon}
\newcommand{\dd}{\,\mathrm{d}}
\begin{document}
\section*{UnsolvedMath ID: 30002700}
\textbf{Source problem code:} catalogue entry
\subsection{Definitions and mathematical statement}
For a smooth Riemannian manifold \((M^n,g)\), a gradient shrinking Ricci soliton consists of a smooth function \(f\colon M\to\R\) and a constant \(\lambda>0\) satisfying
\[
\operatorname{Ric}_g+\nabla_g^2 f=\lambda g.
\]
Here \(\operatorname{Ric}_g\) is the Ricci tensor and \(\nabla_g^2 f\) is the Hessian of \(f\). Let \(\mathcal R_x\colon\Lambda^2T_xM\to\Lambda^2T_xM\) be the self-adjoint curvature operator, with the sign convention that the round sphere has positive curvature operator. Pointwise positive curvature operator means
\[
\langle\mathcal R_x\omega,\omega\rangle_g>0
\quad\text{for every }x\in M\text{ and every }0\ne\omega\in\Lambda^2T_xM.
\]
One interpretation of bounded positive curvature operator is that there exists a constant \(C\in(0,\infty)\), independent of \(x\), such that
\[
0<\langle\mathcal R_x\omega,\omega\rangle_g\le C\lVert\omega\rVert_g^2
\quad\text{for every }x\in M\text{ and }0\ne\omega\in\Lambda^2T_xM.
\]
This condition does not impose a uniform positive lower bound on the curvature operator.

The supplied assertion is that every gradient shrinking Ricci soliton with bounded positive curvature operator is compact. Its wording does not explicitly specify completeness, connectedness, boundary conventions, the dimension range, or whether positivity is uniform. Consequently, the following precise formulation is conditional on interpreting the intended objects as complete, connected manifolds without boundary of dimension at least two, and interpreting bounded positivity as above:

For every integer \(n\ge2\), every complete, connected smooth Riemannian \(n\)-manifold \((M,g)\) without boundary, every \(f\in C^\infty(M)\), and every \(\lambda>0\), if
\[
\operatorname{Ric}_g+\nabla_g^2 f=\lambda g
\]
and there exists \(C\in(0,\infty)\) such that
\[
0<\langle\mathcal R_x\omega,\omega\rangle_g\le C\lVert\omega\rVert_g^2
\quad\text{for all }x\in M,\quad 0\ne\omega\in\Lambda^2T_xM,
\]
must \(M\) be compact?

The added conventions in this conditional formulation require source confirmation; they are not asserted to have been explicitly stated in the supplied report. In particular, completeness cannot be silently inferred from the displayed soliton equation.
\subsection{Short English statement}
The problem asks whether bounded positive curvature operator forces a gradient shrinking Ricci soliton to be compact in every intended dimension. The displayed formulation makes customary completeness and manifold conventions explicit, but these conventions and the meaning of bounded positivity require confirmation. The supplied registry labels the problem open; this draft does not independently verify that status.
\subsection{Sources}
\begin{itemize}
\item Local registry: \texttt{data/unsolvedmath/problems.json}, permanent UnsolvedMath ID \(30002700\), source problem code \texttt{catalogue entry}, title ``Compactness of Positively Curved Shrinking Ricci Solitons.'' The supplied record belongs to the Oberwolfach Reports Open Problems collection. No individual proposer is supplied; the recorded proposed year is 2014.
\item Supplied primary citation: Mini-Workshop: Einstein Metrics, Ricci Solitons and Ricci Flow under Symmetry Assumptions (2015), Oberwolfach Reports, pp.\ 2529--2568, DOI \texttt{10.4171/owr/2014/45}. \url{https://doi.org/10.4171/owr/2014/45}. The literature-source label is ``Oberwolfach Report 45/2014.'' The citation's parenthetical year 2015 and its 2014 report designation are retained as supplied.
\item Supplied literature source: ``Rigidity of Asymptotically Cylindrical Shrinkers.'' \url{https://cvgmt.sns.it/media/doc/paper/2295/CDM-SHRINKERS_3apr.pdf}.
\item Supplied literature source: ``Compact Gradient Shrinking Solitons with Positive Curvature Operator.'' \url{https://arxiv.org/abs/math/0601599}.
\item Literature assessment dated exactly 2026-08-21: the supplied curation labels the problem open and states that compactness is proved in dimensions three and four and under additional curvature, pinching, or volume-growth hypotheses, while the all-dimensional conjecture stated here is not resolved by the available results. This is supplied dated curation derived from the Oberwolfach material, not an independently verified literature assessment; the cited claims and time-sensitive status require independent checking.
\item The registry gives the statement status as \texttt{corrected\_verified} and reports that the clean statement was checked against the official Oberwolfach report and preserves its hypotheses, notation, and open target. These are supplied verification claims, not verification performed for this draft.
\end{itemize}
\end{document}
